\clearpage
\section{Measurement Setup}
\label{sec:setup}
\raggedbottom% this appendix's pages end naturally

\textbf{Instruments.} Every measurement in this datasheet except phase noise was made with a Keysight FieldFox N9918B
handheld analyzer (firmware A.13.25), using its spectrum analyzer and its built-in signal source: the source drives
the unit's input with a CW tone, or sweeps across a window in step with the analyzer as a tracking generator, and
the analyzer reads the unit's output. Phase noise (Section~\ref{sec:phase-noise}) was measured on a Keysight N9042B
UXA signal analyzer with its Phase Noise application. Every setting of the unit is made over the USB serial interface
(Appendix~\ref{sec:commands}), by scripts that run the instruments and the unit together.

\textbf{What ``dBm at the connector'' means.} The test cables between the instruments and the unit lose a few dB of their
own, more at higher frequencies. Each cable's loss (its S21, from 0.5 to 4.5\,GHz for the FR1 cable and 2 to
26.5\,GHz for the FR3 cable) was measured with the FieldFox in network-analyzer mode and is removed from every
reading, frequency by frequency. Every power in this datasheet is therefore referred to the unit's own
connector --- ``dBm at the connector'' --- not to the instrument.

\textbf{Readings.} Two methods are used. Single-tone (CW) measurements, such as compression, sideband suppression,
image rejection and LO leakage, read the tone's power with a peak-search marker on a narrow analyzer span. The unit
and the analyzer run on independent clocks, so a tone can sit slightly off its nominal frequency; the analyzer
re-centers on it before each reading. Swept measurements --- conversion gain and gain flatness, which need the gain
across a whole 500\,MHz window --- use the FieldFox's tracking generator instead: the source and the analyzer sweep
the window together, and the whole trace is read (for gain flatness, 101 points 5\,MHz apart). The swept method was checked against
marker readings of individual tones at the same frequencies. Noise and weak-signal readings use power averaging
(Section~\ref{sec:rx-noise-figure}). Where the input level matters, it was measured directly at the end of the input
cable rather than taken from the source's nominal setting.

\textbf{Configuration.} Unless a section says otherwise, the unit uses its own 10\,MHz TCXO and LO, and all measurements are on Unit~\#5. \textbf{TX side (Sections~\ref{sec:tx-gain-upper}--\ref{sec:spurious-outputs}):} \lmxDriveRule, the TX LO leakage suppressed with the unit's stored calibration (Section~\ref{sec:lo-suppression}), no fan, LMX2820 die at about \condGainMinC--\condCompMaxC\textdegree C (deliberately hot, the worst case). \textbf{RX side:} the same LMX2820 drive rule, with the TX side set up at every point exactly as in normal use (its stored LO-leakage calibration applied, its optimal filter engaged, and both TX-side ports terminated in 50\,\textohm), no fan (the LMX2820 die temperature of each run is given in its section).

\textbf{Measurement automation.} Claude, an AI model made by Anthropic, assisted with writing the measurement scripts
that control the instruments and the unit, and with automating the long frequency and power sweeps, working under
the direction of Pi-Radio engineers. Every number in this datasheet comes from measurements on the real hardware
described above.

\section{Revision History}
\label{sec:revisions}

A datasheet revision \textit{a.b.c} applies to PCB revision \textit{a.b}; the last number counts updates to the
datasheet itself for that PCB revision.

\par\medskip\noindent\begin{minipage}{\linewidth}% caption and table stay on one page
\begin{center}
\captionof{table}{Datasheet revisions.}\label{tab:revisions}
\vspace{0.4em}
\footnotesize
\begin{tabularx}{0.84\linewidth}{|l|l|X|}
\hline
\textbf{Revision} & \textbf{Date} & \textbf{Changes} \\
\hline
1.4.0 & October 2026 & First release of this datasheet, for PCB revision~1.4 with firmware~1.4.1. \\
\hline
\end{tabularx}
\end{center}
\end{minipage}\par
